Small batches may not provide enough queries to occupy the GPU. Splitting
one query's candidates among several blocks exposes more parallel work,
but the blocks must collectively retain the global best candidates.
Each block therefore has capacity for a full shortlist before merging
\cite{johnson_gpu}. Exactness here uses one common score-and-position
order; a distance-only cutoff can choose different tied candidates.

\subsection{Candidate partition and deterministic ordering}

For one query, concatenate the selected IVF lists in probe order into
$n_c$ candidate occurrences. Assume each visited row occurs once and
the index remains fixed. For $\gscan\in\{1,2,4,8\}$, block $(q,j)$ scans
\begin{equation}
 U_j=\left[\left\lfloor\frac{j n_c}{\gscan}\right\rfloor,
              \left\lfloor\frac{(j+1)n_c}{\gscan}\right\rfloor\right),
 \qquad j=0,\ldots,\gscan-1.
 \label{eq:5-partition}
\end{equation}
The intervals denote candidates in this concatenation, not database IDs.
Their sizes differ by at most one, and a long list may cross a boundary.
This balances candidate counts more closely than assigning equal numbers
of lists. Probe-prefix offsets locate each thread's first list; a cursor
advances through subsequent lists. This stage reads primary codes only.

To make local and merged selections agree, equal scores are ordered by
physical position:
\begin{equation}
 K_q(x)=\bigl(\primarydist(q,x),p(x)\bigr),
 \qquad
 (a,p)\prec(b,r)\iff a<b\ \text{or}\ (a=b\ \text{and}\ p<r),
 \label{eq:5-primary-key}
\end{equation}
The position $p(x)$ is unique in the fixed index and does not depend on
the order of atomic appends. This complete key orders every candidate
with a non-NaN score. Final-result ties instead use database IDs.
Changing a partition preserves the selected keys when the computed
scores and this ordering remain the same across paths.

\subsection{Exact local retention and global coverage}

Write $\operatorname{Top}_c(V)$ for the first $\min(c,|V|)$ keys, with
$c=\max\{k,\lceil\alpha k\rceil\}$. All global winners could fall in
one group, so a quota of $c/\gscan$ per group is unsafe. Each group must
retain up to $c$. The following buffer guarantee says that compaction
keeps every local winner still eligible under the shared bound.
The bound $\beta_t$ is introduced in \cref{sec:5-5}; without sharing,
it is a sentinel, a placeholder key later than every valid key.

\begin{resultblock}
\begin{lemma}[Local-buffer retention]
\label{lem:5-buffer}
Let $V_t$ be the candidates scored by a block through chunk $t$, and
$B_t$ its valid buffer after any compaction or filtering at that point.
Assume fixed complete keys, sufficient append capacity and exact
top-$c$ compaction. Let $\beta_t$ be the minimum of the complete shared
bounds observed so far, including the initial top sentinel. Then
\begin{equation}
 \operatorname{Top}_c(V_t)\cap\{x:K_q(x)\preceq\beta_t\}
 \subseteq B_t.
 \label{eq:5-buffer-invariant}
\end{equation}
With sharing disabled, the bound is the sentinel and the buffer contains
the whole prefix top-$c$.
\end{lemma}
\begin{proof}
The empty prefix is immediate. Put
\[
 E_t=\operatorname{Top}_c(V_t)\cap\{x:K_q(x)\preceq\beta_t\}.
\]
On extending $V_t$ to $V_{t+1}$, $\beta_{t+1}\preceq\beta_t$ and
\[
 \operatorname{Top}_c(V_{t+1})\cap V_t\subseteq\operatorname{Top}_c(V_t).
\]
Consequently,
\[
 x\in E_{t+1}\cap V_t\Rightarrow x\in E_t\subseteq B_t.
\]
A new key rejected locally has $c$ predecessors, hence rank $>c$. Capacity admits the remaining eligible arrivals. For any eligible winner $x$,
\[
 \#\{y\in B:K_q(y)\prec K_q(x)\}
 \leq\#\{y\in V_{t+1}:K_q(y)\prec K_q(x)\}<c,
\]
so exact compaction keeps $x$. Filtering removes $K>\beta_{\rm seen}$; the running minimum covers stale reads. These steps preserve the invariant inductively.
\end{proof}
\end{resultblock}
Once every global winner is covered by the retained lists, their exact
top-$c$ can be recovered. The next proposition also permits each merge
to discard everything after its first $c$ entries.

\begin{resultblock}
\begin{proposition}[Coverage and truncated merging]
\label{prop:5-local-coverage}
For fixed keys under a common total order, let $U$ be the visited
universe and $V$ any retained subset. If
$\operatorname{Top}_c(U)\subseteq V\subseteq U$, then
$\operatorname{Top}_c(V)=\operatorname{Top}_c(U)$.
Pairwise sorted merging with truncation to $c$ preserves this result.
In particular, for disjoint groups,
\begin{equation}
 \operatorname{Top}_c\!\left(\bigcup_j U_j\right)
 =\operatorname{Top}_c\!\left(\bigcup_j
                    \operatorname{Top}_c(U_j)\right).
 \label{eq:5-topc-union}
\end{equation}
The statement includes universes smaller than $c$.
\end{proposition}
\begin{proof}
Every global winner is in $V$, and every other member of $V$ follows
the global winners. Selection from $V$ therefore recovers them.
For the independent-group special case, let $x$ belong to the global
top-$c$ and to group $U_j$. Fewer than $c$
candidates precede $x$ globally, so fewer than $c$ precede it within $U_j$.
Consequently $x\in\operatorname{Top}_c(U_j)$, and the local union contains
every global winner. Selecting its first $c$ therefore recovers precisely
the global winners in the same order. If the whole universe has fewer
than $c$ members, every member is retained.

For merging, the first $c$ of two sorted lists are their union's
top-$c$. A discarded candidate already has at least $c$ predecessors
within that union and cannot enter the top-$c$ of any larger union.
Induction over the merge tree establishes the claim.
\end{proof}
\end{resultblock}
\par\medskip\noindent\begin{minipage}{\linewidth}
\captionsetup{hypcap=false}
\centering
% Schematic only: illustrative keys, not experimental observations.
\begin{tikzpicture}[
  font=\small\linespread{1}\selectfont,
  box/.style={draw=ThesisNavy!65,rounded corners=2pt,align=center,
    inner sep=5pt,text width=2.65cm,minimum height=2.05cm},
  merged/.style={box,text width=5.55cm,minimum height=1.05cm,
    fill=ThesisNavy!5},
  arr/.style={-{Latex[length=2mm]},draw=ThesisNavy!80,thick}]
\node[box] (g0) at (0,0) {$\operatorname{Top}_3(U_0)$\\[2pt]
  $\mathbf{(1,8)}$\\$\mathbf{(2,6)}$\\$(2,9)$};
\node[box] (g1) at (3.5,0) {$\operatorname{Top}_3(U_1)$\\[2pt]
  $(4,1)$\\$(5,4)$\\\textit{length 2}};
\node[box] (g2) at (7,0) {$\operatorname{Top}_3(U_2)$\\[2pt]
  $\varnothing$\\\textit{length 0}};
\node[box] (g3) at (10.5,0) {$\operatorname{Top}_3(U_3)$\\[2pt]
  $\mathbf{(2,3)}$\\$(7,2)$\\$(8,5)$};
\node[merged] (m0) at (1.75,-2.5) {Merge and retain first 3\\[2pt]
  $\mathbf{(1,8)},\ \mathbf{(2,6)},\ (2,9)$};
\node[merged] (m1) at (8.75,-2.5) {Merge and retain first 3\\[2pt]
  $\mathbf{(2,3)},\ (7,2),\ (8,5)$};
\node[merged,text width=6.8cm,fill=ThesisNavy!12] (out) at (5.25,-4.3)
  {Global primary top-3\\[2pt]
   $\mathbf{(1,8)},\ \mathbf{(2,3)},\ \mathbf{(2,6)}$};
\draw[arr] (g0.south) -- (m0.north west);
\draw[arr] (g1.south) -- (m0.north east);
\draw[arr] (g2.south) -- (m1.north west);
\draw[arr] (g3.south) -- (m1.north east);
\draw[arr] (m0.south) -- (out.north west);
\draw[arr] (m1.south) -- (out.north east);
\end{tikzpicture}

\captionof{figure}[Local coverage and truncated pairwise merging]{Local coverage and truncated pairwise merging, illustrated with
$c=3$ and $\gscan=4$. Entries are illustrative $(\primarydist,p)$ keys; bold ones are global winners. Short/empty groups contribute their valid prefixes. Merges keep at most $c$ in the same order.}
\label{fig:5-global-selection}
\end{minipage}\par\medskip
Each block holds key pairs in a shared-memory buffer of capacity
$C_{\rm buf}$, a power of two. For $T_{\rm blk}$ scan threads, it reserves
$C_{\rm buf}\geq c+T_{\rm blk}$ entries: room for the retained keys and
one more chunk. Bitonic compaction sorts the buffer and keeps up to $c$
valid keys. Once $c$ exist, the worst retained key becomes the local
threshold $\lambda_j$; until then the threshold is the sentinel.
New candidates enter only when $K\prec\lambda_j$. The $c$ retained
keys justify this threshold even as more candidates arrive.

After compaction, one chunk can add at most one key per thread, so
\[
 |B|+|\Delta V|\leq c+T_{\rm blk}\leq C_{\rm buf}.
\]
This supplies the append-capacity premise of \cref{lem:5-buffer}.
Together with \cref{prop:5-shared-bound,prop:5-local-coverage}, it connects
the buffer implementation to global coverage. Writing every primary
score would require $O(n_c)$ storage and writes; retaining candidates
during the scan emits $O(\gscan c)$ pairs. Register-based selection is
another option \cite{johnson_gpu}.

\subsection{Merge implementation, storage and boundary cases}

For $\gscan>1$, a 256-thread block performs $\log_2\gscan$ merge rounds
per query. A pair with lengths $a,b$ emits $\min(c,a+b)$ keys using
Merge Path slices \cite{gpu_merge_path}. Comparisons follow
\cref{eq:5-primary-key}, with block barriers between rounds.
\Cref{fig:5-global-selection} shows why truncation is safe. The internal
nodes emit at most $c(\gscan-1)$ pairs, in addition to the searches that
partition merge work and the synchronisation overhead.

Two buffers alternate between rounds, with shapes
$[b_{\rm cap},\gscan,c]$ and $[b_{\rm cap},\gscan/2,c]$, where
$b_{\rm cap}$ is the reserved query capacity:
\begin{equation}
 S_{\rm pairs}=8b_{\rm cap}c\left(\gscan+\frac{\gscan}{2}\right)
 \quad\text{bytes},\qquad \gscan>1.
 \label{eq:5-group-workspace}
\end{equation}
Counts add $4b_{\rm cap}\gscan$ bytes and provisioned shared bounds add
$8b_{\rm cap}$ bytes; survivor and other workspaces are separate. These
allocations grow before graph capture and persist for reuse. The merge
block's shared memory holds list lengths rather than all the keys.

Short or empty groups contribute only valid prefixes. Sorting and output
use different padding keys:
\[
 K_{\rm sort\ pad}=(+\infty,\mathrm{INT\_MAX}),\qquad
 K_{\rm output\ pad}=(+\infty,-1).
\]
Valid positions are nonnegative and below $\mathrm{INT\_MAX}$; a negative
output position marks an unused entry. An empty group has count zero,
and $\gscan=1$ writes directly without a merge.

The one-block path, $\gscan=0$, keeps up to $c$ smallest computed scores,
but does not fix which candidates survive an exact tie at the cutoff.
For nonempty $V$, let $h=\min(c,|V|)$ and let $d_h$ be its $h$th score:
\[
 \{x:d_x<d_h\}\subseteq B\subseteq\{x:d_x\leq d_h\},\qquad |B|=h.
\]
Thus every strictly smaller score survives and all retained scores are
at most the cutoff; tied candidate identities may vary. The optional
quota-based approximate selector is disabled in the evaluated builds.
\Cref{sec:7-strict-selection} checks both this distance-order contract
and grouped complete-key ordering by CPU sorting GPU-computed scores.
